% Options for packages loaded elsewhere
% Options for packages loaded elsewhere
\PassOptionsToPackage{unicode}{hyperref}
\PassOptionsToPackage{hyphens}{url}
\PassOptionsToPackage{dvipsnames,svgnames,x11names}{xcolor}
%
\documentclass[
  letterpaper,
  DIV=11,
  numbers=noendperiod]{scrartcl}
\usepackage{xcolor}
\usepackage{amsmath,amssymb}
\setcounter{secnumdepth}{-\maxdimen} % remove section numbering
\usepackage{iftex}
\ifPDFTeX
  \usepackage[T1]{fontenc}
  \usepackage[utf8]{inputenc}
  \usepackage{textcomp} % provide euro and other symbols
\else % if luatex or xetex
  \usepackage{unicode-math} % this also loads fontspec
  \defaultfontfeatures{Scale=MatchLowercase}
  \defaultfontfeatures[\rmfamily]{Ligatures=TeX,Scale=1}
\fi
\usepackage{lmodern}
\ifPDFTeX\else
  % xetex/luatex font selection
\fi
% Use upquote if available, for straight quotes in verbatim environments
\IfFileExists{upquote.sty}{\usepackage{upquote}}{}
\IfFileExists{microtype.sty}{% use microtype if available
  \usepackage[]{microtype}
  \UseMicrotypeSet[protrusion]{basicmath} % disable protrusion for tt fonts
}{}
\makeatletter
\@ifundefined{KOMAClassName}{% if non-KOMA class
  \IfFileExists{parskip.sty}{%
    \usepackage{parskip}
  }{% else
    \setlength{\parindent}{0pt}
    \setlength{\parskip}{6pt plus 2pt minus 1pt}}
}{% if KOMA class
  \KOMAoptions{parskip=half}}
\makeatother
% Make \paragraph and \subparagraph free-standing
\makeatletter
\ifx\paragraph\undefined\else
  \let\oldparagraph\paragraph
  \renewcommand{\paragraph}{
    \@ifstar
      \xxxParagraphStar
      \xxxParagraphNoStar
  }
  \newcommand{\xxxParagraphStar}[1]{\oldparagraph*{#1}\mbox{}}
  \newcommand{\xxxParagraphNoStar}[1]{\oldparagraph{#1}\mbox{}}
\fi
\ifx\subparagraph\undefined\else
  \let\oldsubparagraph\subparagraph
  \renewcommand{\subparagraph}{
    \@ifstar
      \xxxSubParagraphStar
      \xxxSubParagraphNoStar
  }
  \newcommand{\xxxSubParagraphStar}[1]{\oldsubparagraph*{#1}\mbox{}}
  \newcommand{\xxxSubParagraphNoStar}[1]{\oldsubparagraph{#1}\mbox{}}
\fi
\makeatother

\usepackage{color}
\usepackage{fancyvrb}
\newcommand{\VerbBar}{|}
\newcommand{\VERB}{\Verb[commandchars=\\\{\}]}
\DefineVerbatimEnvironment{Highlighting}{Verbatim}{commandchars=\\\{\}}
% Add ',fontsize=\small' for more characters per line
\usepackage{framed}
\definecolor{shadecolor}{RGB}{241,243,245}
\newenvironment{Shaded}{\begin{snugshade}}{\end{snugshade}}
\newcommand{\AlertTok}[1]{\textcolor[rgb]{0.68,0.00,0.00}{#1}}
\newcommand{\AnnotationTok}[1]{\textcolor[rgb]{0.37,0.37,0.37}{#1}}
\newcommand{\AttributeTok}[1]{\textcolor[rgb]{0.40,0.45,0.13}{#1}}
\newcommand{\BaseNTok}[1]{\textcolor[rgb]{0.68,0.00,0.00}{#1}}
\newcommand{\BuiltInTok}[1]{\textcolor[rgb]{0.00,0.23,0.31}{#1}}
\newcommand{\CharTok}[1]{\textcolor[rgb]{0.13,0.47,0.30}{#1}}
\newcommand{\CommentTok}[1]{\textcolor[rgb]{0.37,0.37,0.37}{#1}}
\newcommand{\CommentVarTok}[1]{\textcolor[rgb]{0.37,0.37,0.37}{\textit{#1}}}
\newcommand{\ConstantTok}[1]{\textcolor[rgb]{0.56,0.35,0.01}{#1}}
\newcommand{\ControlFlowTok}[1]{\textcolor[rgb]{0.00,0.23,0.31}{\textbf{#1}}}
\newcommand{\DataTypeTok}[1]{\textcolor[rgb]{0.68,0.00,0.00}{#1}}
\newcommand{\DecValTok}[1]{\textcolor[rgb]{0.68,0.00,0.00}{#1}}
\newcommand{\DocumentationTok}[1]{\textcolor[rgb]{0.37,0.37,0.37}{\textit{#1}}}
\newcommand{\ErrorTok}[1]{\textcolor[rgb]{0.68,0.00,0.00}{#1}}
\newcommand{\ExtensionTok}[1]{\textcolor[rgb]{0.00,0.23,0.31}{#1}}
\newcommand{\FloatTok}[1]{\textcolor[rgb]{0.68,0.00,0.00}{#1}}
\newcommand{\FunctionTok}[1]{\textcolor[rgb]{0.28,0.35,0.67}{#1}}
\newcommand{\ImportTok}[1]{\textcolor[rgb]{0.00,0.46,0.62}{#1}}
\newcommand{\InformationTok}[1]{\textcolor[rgb]{0.37,0.37,0.37}{#1}}
\newcommand{\KeywordTok}[1]{\textcolor[rgb]{0.00,0.23,0.31}{\textbf{#1}}}
\newcommand{\NormalTok}[1]{\textcolor[rgb]{0.00,0.23,0.31}{#1}}
\newcommand{\OperatorTok}[1]{\textcolor[rgb]{0.37,0.37,0.37}{#1}}
\newcommand{\OtherTok}[1]{\textcolor[rgb]{0.00,0.23,0.31}{#1}}
\newcommand{\PreprocessorTok}[1]{\textcolor[rgb]{0.68,0.00,0.00}{#1}}
\newcommand{\RegionMarkerTok}[1]{\textcolor[rgb]{0.00,0.23,0.31}{#1}}
\newcommand{\SpecialCharTok}[1]{\textcolor[rgb]{0.37,0.37,0.37}{#1}}
\newcommand{\SpecialStringTok}[1]{\textcolor[rgb]{0.13,0.47,0.30}{#1}}
\newcommand{\StringTok}[1]{\textcolor[rgb]{0.13,0.47,0.30}{#1}}
\newcommand{\VariableTok}[1]{\textcolor[rgb]{0.07,0.07,0.07}{#1}}
\newcommand{\VerbatimStringTok}[1]{\textcolor[rgb]{0.13,0.47,0.30}{#1}}
\newcommand{\WarningTok}[1]{\textcolor[rgb]{0.37,0.37,0.37}{\textit{#1}}}

\usepackage{longtable,booktabs,array}
\newcounter{none} % for unnumbered tables
\usepackage{calc} % for calculating minipage widths
% Correct order of tables after \paragraph or \subparagraph
\usepackage{etoolbox}
\makeatletter
\patchcmd\longtable{\par}{\if@noskipsec\mbox{}\fi\par}{}{}
\makeatother
% Allow footnotes in longtable head/foot
\IfFileExists{footnotehyper.sty}{\usepackage{footnotehyper}}{\usepackage{footnote}}
\makesavenoteenv{longtable}
\usepackage{graphicx}
\makeatletter
\newsavebox\pandoc@box
\newcommand*\pandocbounded[1]{% scales image to fit in text height/width
  \sbox\pandoc@box{#1}%
  \Gscale@div\@tempa{\textheight}{\dimexpr\ht\pandoc@box+\dp\pandoc@box\relax}%
  \Gscale@div\@tempb{\linewidth}{\wd\pandoc@box}%
  \ifdim\@tempb\p@<\@tempa\p@\let\@tempa\@tempb\fi% select the smaller of both
  \ifdim\@tempa\p@<\p@\scalebox{\@tempa}{\usebox\pandoc@box}%
  \else\usebox{\pandoc@box}%
  \fi%
}
% Set default figure placement to htbp
\def\fps@figure{htbp}
\makeatother





\setlength{\emergencystretch}{3em} % prevent overfull lines

\providecommand{\tightlist}{%
  \setlength{\itemsep}{0pt}\setlength{\parskip}{0pt}}



 


\KOMAoption{captions}{tableheading}
\makeatletter
\@ifpackageloaded{caption}{}{\usepackage{caption}}
\AtBeginDocument{%
\ifdefined\contentsname
  \renewcommand*\contentsname{Table of contents}
\else
  \newcommand\contentsname{Table of contents}
\fi
\ifdefined\listfigurename
  \renewcommand*\listfigurename{List of Figures}
\else
  \newcommand\listfigurename{List of Figures}
\fi
\ifdefined\listtablename
  \renewcommand*\listtablename{List of Tables}
\else
  \newcommand\listtablename{List of Tables}
\fi
\ifdefined\figurename
  \renewcommand*\figurename{Figure}
\else
  \newcommand\figurename{Figure}
\fi
\ifdefined\tablename
  \renewcommand*\tablename{Table}
\else
  \newcommand\tablename{Table}
\fi
}
\@ifpackageloaded{float}{}{\usepackage{float}}
\floatstyle{ruled}
\@ifundefined{c@chapter}{\newfloat{codelisting}{h}{lop}}{\newfloat{codelisting}{h}{lop}[chapter]}
\floatname{codelisting}{Listing}
\newcommand*\listoflistings{\listof{codelisting}{List of Listings}}
\makeatother
\makeatletter
\makeatother
\makeatletter
\@ifpackageloaded{caption}{}{\usepackage{caption}}
\@ifpackageloaded{subcaption}{}{\usepackage{subcaption}}
\makeatother
\usepackage{bookmark}
\IfFileExists{xurl.sty}{\usepackage{xurl}}{} % add URL line breaks if available
\urlstyle{same}
\hypersetup{
  pdftitle={Interpreting Descriptive and Exploratory Outputs: Plots and Summaries.},
  pdfauthor={Rose Megirian},
  colorlinks=true,
  linkcolor={blue},
  filecolor={Maroon},
  citecolor={Blue},
  urlcolor={Blue},
  pdfcreator={LaTeX via pandoc}}


\title{Interpreting Descriptive and Exploratory Outputs: Plots and
Summaries.}
\author{Rose Megirian}
\date{}
\begin{document}
\maketitle

\renewcommand*\contentsname{Table of contents}
{
\setcounter{tocdepth}{3}
\tableofcontents
}

\begin{figure}

\begin{minipage}[t]{0.33\linewidth}

\begin{figure}[H]

{\centering \pandocbounded{\includegraphics[keepaspectratio]{descriptive_and_exploratory_outputs_files/figure-pdf/header-plotsandfigures-1.pdf}}

}

\subcaption{\hyperref[sec-histograms]{Histograms}}

\end{figure}%

\end{minipage}%
%
\begin{minipage}[t]{0.33\linewidth}

\begin{figure}[H]

{\centering \pandocbounded{\includegraphics[keepaspectratio]{descriptive_and_exploratory_outputs_files/figure-pdf/header-plotsandfigures-2.pdf}}

}

\subcaption{\hyperref[sec-boxplots]{Box Plots}}

\end{figure}%

\end{minipage}%
%
\begin{minipage}[t]{0.33\linewidth}

\begin{figure}[H]

{\centering \pandocbounded{\includegraphics[keepaspectratio]{descriptive_and_exploratory_outputs_files/figure-pdf/header-plotsandfigures-3.pdf}}

}

\subcaption{\hyperref[sec-scatterplots]{Scatter Plots}}

\end{figure}%

\end{minipage}%

\end{figure}%

Descriptive and exploratory analysis summarises and describes the
characteristics of a dataset, providing insight into distributions,
central tendency, spread, and relationships between variables. It serves
two closely related purposes: as a standalone way of characterising and
communicating what the data shows, and as an essential first step before
formal modelling or hypothesis testing. Used in this exploratory
capacity, descriptive analysis helps reveal the shape and structure of
your data, identify potential issues, and inform decisions about how
best to proceed with formal analysis.

This guide covers the most common outputs from descriptive and
exploratory analysis that you are likely to encounter, explaining what
they represent, what different patterns indicate, and what types of data
each is appropriate for. It is organised into two sections: plots and
figures, which covers common visualisations used to explore and describe
data, and numerical summaries, which covers the key statistics used to
characterise and communicate the properties of a dataset.

\section{Plots and Figures}\label{sec-plotsandfigures}

Visualisations are one of the most powerful tools in data analysis,
making patterns, distributions, and relationships visible in ways that
raw numbers alone often cannot. Plots like histograms, boxplots, and
scatter plots are particularly valuable during exploratory data
analysis, helping to reveal the shape and spread of data, identify
outliers, and assess relationships between variables. This guide covers
the most common descriptive plot types you are likely to encounter,
explaining what they represent, what different patterns indicate, and
what types of data each is appropriate for.

\subsection{Histograms}\label{sec-histograms}

\begin{Shaded}
\begin{Highlighting}[]
\NormalTok{//| echo: false}

\NormalTok{viewof dataset = Inputs.select(}
\NormalTok{  [}
\NormalTok{    "Normal Distribution", "Skewed Distribution", "Bimodal Distribution",}
\NormalTok{    "Uniform Distribution", "Exponential Distribution"}
\NormalTok{  ],}
\NormalTok{  \{label: "Distribution", value: "Normal Distribution"\}}
\NormalTok{)}

\NormalTok{viewof bin\_width = Inputs.range([1, 30], \{}
\NormalTok{  label: "Bin width",}
\NormalTok{  step: 1,}
\NormalTok{  value: 6}
\NormalTok{\})}
\end{Highlighting}
\end{Shaded}

\begin{Shaded}
\begin{Highlighting}[]
\NormalTok{//| echo: false}

\NormalTok{html\textasciigrave{}\textless{}div style="margin{-}top: 2em;"\textgreater{}\textless{}/div\textgreater{}\textasciigrave{}}
\end{Highlighting}
\end{Shaded}

\begin{Shaded}
\begin{Highlighting}[]
\NormalTok{//| echo: false}

\NormalTok{all\_data = transpose(hist\_data)}

\NormalTok{filtered = all\_data.filter(d =\textgreater{} d.dataset === dataset)}

\NormalTok{annotation = (\{}
\NormalTok{  "Normal Distribution": \textasciigrave{}Values cluster symmetrically around a central point. }
\NormalTok{    This bell{-}shaped curve is the most commonly assumed distribution in statistical }
\NormalTok{    methods. Try adjusting bin width to see how it affects the apparent shape.\textasciigrave{},}

\NormalTok{  "Skewed Distribution": \textasciigrave{}This distribution has a long tail to the right, meaning most }
\NormalTok{    values are low but a few are very high. This is called right skew, or positive skew. }
\NormalTok{    Skewed data can violate the normality assumptions of some statistical methods.\textasciigrave{},}

\NormalTok{  "Bimodal Distribution": \textasciigrave{}Two distinct peaks suggest two subgroups within the data. }
\NormalTok{    This is called a bimodal distribution. Seeing two humps in a histogram is often }
\NormalTok{    a signal to investigate whether the data should be analysed as separate groups.\textasciigrave{},}

\NormalTok{  "Uniform Distribution": \textasciigrave{}Values are spread roughly evenly across the range with no }
\NormalTok{    clear peak. This uniform distribution is relatively rare in practice but useful }
\NormalTok{    for understanding how histograms represent shape.\textasciigrave{},}

\NormalTok{  "Exponential Distribution": \textasciigrave{}A sharp peak near zero with a long right tail. Common }
\NormalTok{    in data representing waiting times or survival times. Heavily skewed data like }
\NormalTok{    this often benefits from a log transformation before analysis.\textasciigrave{}}

\NormalTok{\})[dataset] || ""}

\NormalTok{Plot.plot(\{}
\NormalTok{  style: \{background: "transparent", fontFamily: "inherit"\},}
\NormalTok{  marginBottom: 50,}
\NormalTok{  x: \{label: "Value"\},}
\NormalTok{  y: \{label: "Frequency"\},}
\NormalTok{  marks: [}
\NormalTok{    Plot.rectY(}
\NormalTok{      filtered,}
\NormalTok{      Plot.binX(}
\NormalTok{        \{y: "count"\},}
\NormalTok{        \{}
\NormalTok{          x: d =\textgreater{} +d.value,}
\NormalTok{          thresholds: d3.range(}
\NormalTok{            d3.min(filtered, d =\textgreater{} +d.value),}
\NormalTok{            d3.max(filtered, d =\textgreater{} +d.value) + bin\_width,}
\NormalTok{            bin\_width}
\NormalTok{          ),}
\NormalTok{          fill: "\#F2F2F2",}
\NormalTok{          stroke: "\#414042",}
\NormalTok{          strokeWidth: 0.8}
\NormalTok{        \}}
\NormalTok{      )}
\NormalTok{    ),}
\NormalTok{    Plot.ruleY([0])}
\NormalTok{  ]}
\NormalTok{\})}
\end{Highlighting}
\end{Shaded}

\begin{Shaded}
\begin{Highlighting}[]
\NormalTok{//| echo: false}

\NormalTok{html\textasciigrave{}\textless{}div style="font{-}size: 0.9em; color: \#555; margin{-}top: .5em; }
\NormalTok{  padding: 0.75em 1em; background: \#f9f9f9; border{-}left: 3px solid \#414042;"\textgreater{}}
\NormalTok{  $\{annotation\}}
\NormalTok{\textless{}/div\textgreater{}\textasciigrave{}}
\end{Highlighting}
\end{Shaded}

A histogram displays the distribution of a continuous variable by
grouping values into intervals known as bins, and showing how many
observations fall within each bin as a bar. The height of each bar
represents the frequency of values in that range, making it easy to see
at a glance where values cluster, how spread out they are, and whether
any unusual patterns are present.

A key feature of a histogram is that the bars touch. This reflects the
fact that the underlying data is continuous, meaning there are no gaps
between possible values and therefore no gaps between bins. This
distinguishes a histogram from a bar chart, which is used for discrete
or categorical data where each bar represents a separate group rather
than a range of a continuous scale. If your data consists of counts,
ordinal scores, or named categories, a bar chart is the more appropriate
choice.

Histograms are commonly used during exploratory data analysis to
understand the distribution of a variable before any formal modelling
takes place. They help reveal how spread out values are, whether they
cluster around a central point, whether there are unusually high or low
values, and whether the distribution is symmetrical or skewed. This kind
of early understanding of your data is important for informing decisions
about which analytical methods are likely to be appropriate.

Histograms are also widely used during and after model fitting to
examine residuals, the differences between observed values and the
values predicted by the model. Most common analytical methods require
that residuals are approximately normally distributed, and plotting a
histogram of residuals is a straightforward way to assess whether this
holds. It is worth noting that it is the residuals rather than the raw
data that need to follow an approximately normal distribution. Raw data
can be skewed or non-normal and the residuals can still be approximately
normal after a model is fitted. However, heavily non-normal raw data may
be an early indication that this is worth checking carefully once a
model has been fitted.

\subsubsection{Reading histograms}\label{reading-histograms}

Look at where the bars are tallest, as this is where most values fall. A
symmetrical bell-shaped curve suggests an approximately normal
distribution, where values cluster evenly around a central point with
the mean, median and mode all close together
(Figure~\ref{fig-hist-normal}). A longer tail on one side indicates
skewness. A right-skewed distribution has a tail extending to the right,
meaning most values are lower with fewer very high values pulling the
mean upward. A left-skewed distribution has a tail extending to the
left, meaning most values are higher with fewer very low values pulling
the mean downward (Figure~\ref{fig-hist-skewed}). Two distinct peaks
suggest the data may contain two subgroups with different underlying
characteristics, known as a bimodal distribution, which is worth
investigating further before proceeding with analysis
(Figure~\ref{fig-hist-bimodal}). Bars sitting far from the main cluster
may indicate outliers worth examining.

\paragraph{Normal example}

\begin{figure}

\centering{

\pandocbounded{\includegraphics[keepaspectratio]{descriptive_and_exploratory_outputs_files/figure-pdf/fig-hist-normal-1.pdf}}

}

\caption{\label{fig-hist-normal}Histogram of an approximately normal
distribution, showing the mean, median and mode clustering around the
centre of a symmetrical bell-shaped curve. This is the most commonly
assumed distribution shape underlying statistical methods, making it an
important pattern to recognise in your data and residuals.}

\end{figure}%

\paragraph{Skewed example}

\begin{figure}

\centering{

\pandocbounded{\includegraphics[keepaspectratio]{descriptive_and_exploratory_outputs_files/figure-pdf/fig-hist-skewed-1.pdf}}

}

\caption{\label{fig-hist-skewed}A histogram of a skewed distribution,
where the longer tail pulls the mean away from the median and mode. This
separation of the three measures is a key indicator of skewness and is
important for determining which analytical methods are appropriate.}

\end{figure}%

\paragraph{Bimodal example}

\begin{figure}

\centering{

\pandocbounded{\includegraphics[keepaspectratio]{descriptive_and_exploratory_outputs_files/figure-pdf/fig-hist-bimodal-1.pdf}}

}

\caption{\label{fig-hist-bimodal}Histogram of a bimodal distribution
showing two distinct peaks, suggesting the data may contain two
subgroups with different underlying characteristics. Applying an
analysis that assumes a single trend across all the data, such as
comparing a single overall mean, could obscure these differences and
lead to misleading conclusions.}

\end{figure}%

\subsection{Boxplots}\label{sec-boxplots}

\begin{figure}

\centering{

\pandocbounded{\includegraphics[keepaspectratio]{descriptive_and_exploratory_outputs_files/figure-pdf/fig-box-anatomy-1.pdf}}

}

\caption{\label{fig-box-anatomy}A boxplot provides a concise visual
summary of a variable's distribution. The central box spans the
interquartile range (IQR), from the first quartile (Q1) to the third
quartile (Q3), representing the middle 50\% of values. A line inside the
box marks the median. The whiskers extend to the furthest values that
fall within 1.5 times the IQR below Q1 and above Q3. values beyond this
range are considered potential outliers and are plotted as individual
points, indicating they may require further investigation. In this
example, the relatively balanced spacing of the box and whiskers
suggests that the data may follow a normal distribution.}

\end{figure}%

A boxplot, or box-and-whisker plot, summarises the distribution of a
continuous variable by displaying key statistics including the median,
quartiles, and potential outliers in a compact visual form. Boxplots are
particularly useful in exploratory data analysis for understanding the
spread and shape of your data, and for comparing distributions across
groups.

Because the boxplot is built around the median and quartiles rather than
the mean, it is not strongly influenced by extreme values. A single very
high or very low observation can pull the mean substantially away from
the centre of the data, but has little effect on the median or the box,
making boxplots a reliable summary of your data's distribution even when
outliers are present.

It is worth noting that the 1.5 times IQR rule used to define the
whiskers and flag potential outliers is a convention rather than a
formal statistical test. Points plotted beyond the whiskers are not
necessarily problematic but are worth examining in the context of your
data. It is also worth being aware that boxplots do not convey sample
size, meaning a box based on ten observations looks the same as one
based on ten thousand.

\subsubsection{Reading boxplots}\label{reading-boxplots}

\hfill\break

\begin{Shaded}
\begin{Highlighting}[]
\NormalTok{//| echo: false}

\NormalTok{viewof box\_dataset = Inputs.select(}
\NormalTok{  [}
\NormalTok{    "Normal Distribution", "Skewed Distribution", "Bimodal Distribution",}
\NormalTok{    "Uniform Distribution", "Exponential Distribution"}
\NormalTok{  ],}
\NormalTok{  \{label: "Distribution", value: "Normal Distribution"\}}
\NormalTok{)}
\end{Highlighting}
\end{Shaded}

\begin{Shaded}
\begin{Highlighting}[]
\NormalTok{//| echo: false}

\NormalTok{html\textasciigrave{}\textless{}div style="margin{-}top: 2em;"\textgreater{}\textless{}/div\textgreater{}\textasciigrave{}}
\end{Highlighting}
\end{Shaded}

\begin{Shaded}
\begin{Highlighting}[]
\NormalTok{//| echo: false}

\NormalTok{box\_all\_data = transpose(box\_data)}

\NormalTok{box\_filtered = box\_all\_data.filter(d =\textgreater{} d.dataset === box\_dataset)}

\NormalTok{box\_values = box\_filtered.map(d =\textgreater{} +d.value)}

\NormalTok{// Compute summary stats}
\NormalTok{box\_sorted = [...box\_values].sort((a, b) =\textgreater{} a {-} b)}

\NormalTok{box\_q1 = d3.quantile(box\_sorted, 0.25)}
\NormalTok{box\_q3 = d3.quantile(box\_sorted, 0.75)}
\NormalTok{box\_median = d3.quantile(box\_sorted, 0.5)}
\NormalTok{box\_iqr = box\_q3 {-} box\_q1}
\NormalTok{box\_lower\_fence = box\_q1 {-} 1.5 * box\_iqr}
\NormalTok{box\_upper\_fence = box\_q3 + 1.5 * box\_iqr}
\NormalTok{box\_lower\_whisker = d3.min(box\_values.filter(d =\textgreater{} d \textgreater{}= box\_lower\_fence))}
\NormalTok{box\_upper\_whisker = d3.max(box\_values.filter(d =\textgreater{} d \textless{}= box\_upper\_fence))}
\NormalTok{box\_outliers = box\_values.filter(d =\textgreater{} d \textless{} box\_lower\_fence || d \textgreater{} box\_upper\_fence)}

\NormalTok{box\_annotation = (\{}
\NormalTok{  "Normal Distribution": \textasciigrave{}The median sits centrally within the box and the whiskers }
\NormalTok{    are roughly equal in length, indicating a roughly symmetric distribution with }
\NormalTok{    values spread evenly around the centre. This is consistent with a normal }
\NormalTok{    distribution. There are few or no outliers.\textasciigrave{},}

\NormalTok{  "Skewed Distribution": \textasciigrave{}The median is pulled toward the left of the box and the }
\NormalTok{    right whisker is considerably longer than the left, indicating right skew. A }
\NormalTok{    small number of very high values extend the upper whisker and appear as outliers. }
\NormalTok{    Compare this with the histogram view — the long right tail is visible in both }
\NormalTok{    plot types but manifests differently.\textasciigrave{},}

\NormalTok{  "Bimodal Distribution": \textasciigrave{}A bimodal distribution is one of the harder patterns to }
\NormalTok{    detect in a boxplot. The box and whiskers may appear relatively wide, reflecting }
\NormalTok{    the spread between two clusters, but the two peaks are not visible the way they }
\NormalTok{    are in a histogram. This is an important limitation of boxplots — they summarise }
\NormalTok{    distribution shape but can obscure underlying structure. A histogram or density }
\NormalTok{    plot is better for detecting bimodality.\textasciigrave{},}

\NormalTok{  "Uniform Distribution": \textasciigrave{}The median sits near the centre of the box and both }
\NormalTok{    whiskers are long and roughly equal, reflecting the even spread of values across }
\NormalTok{    the full range. There are few outliers because no values fall unusually far from }
\NormalTok{    the rest — everything is equally spread.\textasciigrave{},}

\NormalTok{  "Exponential Distribution": \textasciigrave{}The median sits close to the left edge of the box and }
\NormalTok{    the right whisker is much longer than the left, with many high{-}value outliers }
\NormalTok{    visible beyond the upper whisker. This strongly right{-}skewed pattern is typical }
\NormalTok{    of exponential data and is a signal that standard analytical methods assuming }
\NormalTok{    normality may not be appropriate without transformation.\textasciigrave{}}

\NormalTok{\})[box\_dataset] || ""}

\NormalTok{// Combine whisker and box data for Plot}
\NormalTok{box\_summary = [\{ }
\NormalTok{  x1: box\_lower\_whisker, }
\NormalTok{  x2: box\_upper\_whisker,}
\NormalTok{  q1: box\_q1,}
\NormalTok{  q3: box\_q3,}
\NormalTok{  median: box\_median}
\NormalTok{\}]}

\NormalTok{Plot.plot(\{}
\NormalTok{  style: \{background: "transparent", fontFamily: "inherit"\},}
\NormalTok{  marginBottom: 40,}
\NormalTok{  marginLeft: 40,}
\NormalTok{  height: 160,}
\NormalTok{  x: \{label: "Value"\},}
\NormalTok{  y: \{axis: null\},}
\NormalTok{  marks: [}
\NormalTok{    // Whisker line}
\NormalTok{    Plot.ruleX(box\_summary, \{}
\NormalTok{      x1: d =\textgreater{} d.x1,}
\NormalTok{      x2: d =\textgreater{} d.x2,}
\NormalTok{      y: 0,}
\NormalTok{      stroke: "\#414042",}
\NormalTok{      strokeWidth: 1.5}
\NormalTok{    \}),}
\NormalTok{    // Box}
\NormalTok{    Plot.rect(box\_summary, \{}
\NormalTok{      x1: d =\textgreater{} d.q1,}
\NormalTok{      x2: d =\textgreater{} d.q3,}
\NormalTok{      y1: {-}0.4,}
\NormalTok{      y2: 0.4,}
\NormalTok{      fill: "\#F2F2F2",}
\NormalTok{      stroke: "\#414042",}
\NormalTok{      strokeWidth: 1.5}
\NormalTok{    \}),}
\NormalTok{    // Median line}
\NormalTok{    Plot.ruleX(box\_summary, \{}
\NormalTok{      x: d =\textgreater{} d.median,}
\NormalTok{      y1: {-}0.4,}
\NormalTok{      y2: 0.4,}
\NormalTok{      stroke: "\#414042",}
\NormalTok{      strokeWidth: 2}
\NormalTok{    \}),}
\NormalTok{    // Outliers}
\NormalTok{    Plot.dot(box\_outliers.map(v =\textgreater{} (\{v\})), \{}
\NormalTok{      x: d =\textgreater{} d.v,}
\NormalTok{      y: 0,}
\NormalTok{      fill: "\#ec8525",}
\NormalTok{      r: 3}
\NormalTok{    \}),}
\NormalTok{    // Lower whisker (from lower whisker end to Q1)}
\NormalTok{Plot.link(box\_summary, \{}
\NormalTok{  x1: d =\textgreater{} d.x1,}
\NormalTok{  x2: d =\textgreater{} d.q1,}
\NormalTok{  y1: 0,}
\NormalTok{  y2: 0,}
\NormalTok{  stroke: "\#414042",}
\NormalTok{  strokeWidth: 1.5}
\NormalTok{\}),}
\NormalTok{// Upper whisker (from Q3 to upper whisker end)}
\NormalTok{Plot.link(box\_summary, \{}
\NormalTok{  x1: d =\textgreater{} d.q3,}
\NormalTok{  x2: d =\textgreater{} d.x2,}
\NormalTok{  y1: 0,}
\NormalTok{  y2: 0,}
\NormalTok{  stroke: "\#414042",}
\NormalTok{  strokeWidth: 1.5}
\NormalTok{\})}
\NormalTok{  ]}
\NormalTok{\})}
\end{Highlighting}
\end{Shaded}

\begin{Shaded}
\begin{Highlighting}[]
\NormalTok{//| echo: false}
\NormalTok{html\textasciigrave{}\textless{}div style="font{-}size: 0.9em; color: \#555; margin{-}top: 0.5em; }
\NormalTok{  padding: 0.75em 1em; background: \#f9f9f9; border{-}left: 3px solid \#414042;"\textgreater{}}
\NormalTok{  $\{box\_annotation\}}
\NormalTok{\textless{}/div\textgreater{}\textasciigrave{}}
\end{Highlighting}
\end{Shaded}

When reading a boxplot, start with the median line. Its position within
the box indicates whether the data is roughly symmetric or skewed. A
median sitting centrally suggests an even distribution on both sides,
while a median closer to one end of the box indicates that values are
more concentrated on that side. The relative length of the whiskers
tells a similar story; a longer whisker on one side suggests the data
extends further in that direction.

A wide box indicates high variability in the middle 50\% of values,
while a narrow box suggests values are tightly clustered. The
interactive examples above allow you to see how these features change
across different distribution shapes, and how the same patterns
identified in histograms appear differently when summarised as a
boxplot.

\subsubsection{Comparing groups}\label{comparing-groups}

\begin{figure}

\centering{

\pandocbounded{\includegraphics[keepaspectratio]{descriptive_and_exploratory_outputs_files/figure-pdf/fig-box-comparing-groups-1.pdf}}

}

\caption{\label{fig-box-comparing-groups}Boxplots comparing yield
distributions across three crop varieties. Differences in the position
of the median line indicate differences in central tendency between
groups, while differences in box width and whisker length indicate
differences in spread. Overlapping boxes suggest the groups may not
differ meaningfully, while clearly separated boxes suggest real
differences worth investigating further.}

\end{figure}%

Where boxplots are most powerful is in side-by-side comparison
(Figure~\ref{fig-box-comparing-groups}). Differences in median position
indicate differences in central tendency between groups. Differences in
box width and whisker length indicate differences in spread. Overlapping
boxes suggest groups may not differ meaningfully, while clearly
separated boxes suggest real differences worth investigating further.

\subsection{Scatter Plots}\label{sec-scatterplots}

A scatter plot displays the relationship between two numeric variables
by plotting individual observations as points, with one variable on each
axis. It is most appropriate when both variables are measured on a
continuous scale, making it easy to see at a glance whether a
relationship exists, what direction it takes, and how strong it appears
to be.

Scatter plots are particularly valuable during exploratory data analysis
as an early step before formal modelling. Seeing the shape of the
relationship between two variables helps inform which analytical methods
are likely to be appropriate.\\
\strut \\

\begin{figure}

\begin{minipage}[t]{0.50\linewidth}

\begin{figure}[H]

\centering{

\pandocbounded{\includegraphics[keepaspectratio]{descriptive_and_exploratory_outputs_files/figure-pdf/fig-scatter-linear-1.pdf}}

}

\caption{\label{fig-scatter-linear}A scatter plot showing a strong
positive linear relationship between two variables. Points cluster
tightly around the fitted line, indicating that values of one variable
are highly predictable from the other. The shaded band represents the
confidence interval around the fitted line, reflecting uncertainty in
the estimated relationship.}

\end{figure}%

\end{minipage}%
%
\begin{minipage}[t]{0.50\linewidth}

\begin{figure}[H]

\centering{

\pandocbounded{\includegraphics[keepaspectratio]{descriptive_and_exploratory_outputs_files/figure-pdf/fig-scatter-curved-1.pdf}}

}

\caption{\label{fig-scatter-curved}A scatter plot showing a non-linear
relationship between two variables. The curved fitted line indicates
that the rate of change is not constant across the range of values,
meaning a straight line would poorly describe the relationship.}

\end{figure}%

\end{minipage}%

\end{figure}%

\subsubsection{Reading scatter plots}\label{reading-scatter-plots}

\hfill\break

\begin{figure}

\centering{

\pandocbounded{\includegraphics[keepaspectratio]{descriptive_and_exploratory_outputs_files/figure-pdf/fig-scatter-groups-1.pdf}}

}

\caption{\label{fig-scatter-groups}A scatter plot showing two distinct
groups with different relationships between the same two variables. The
steeper cluster of points indicates a more rapid rate of change, while
the shallower cluster indicates a slower rate. When groups show
different relationships like this, analysing them together as a single
trend can be misleading, and the grouping variable should be accounted
for in any subsequent analysis.}

\end{figure}%

Start with the overall pattern. Points forming an upward trend indicate
a positive relationship, where higher values of one variable tend to
coincide with higher values of the other. A downward trend indicates a
negative relationship. A random scatter with no clear direction suggests
little or no relationship between the variables.

The closeness of points to an imaginary or fitted line reflects the
strength of the relationship. Points clustering tightly around a clear
trend indicate a strong relationship, while a more dispersed pattern
suggests a weaker one. It is also worth looking at whether the
relationship is linear, following a straight line, or non-linear,
following a curve, as this affects which analytical methods are
appropriate (Figure~\ref{fig-scatter-linear},
Figure~\ref{fig-scatter-curved}).

Where points fall into distinct clusters or groups, analysing them as a
single trend can be misleading. If a grouping variable is present in
your data, it is worth plotting groups separately or using colour to
distinguish them, as different groups may show quite different
relationships (Figure~\ref{fig-scatter-groups}).

With large datasets, points can overlap heavily, making it difficult to
see the true density and distribution of observations. If a scatter plot
looks unusually sparse or clustered in a way that seems inconsistent
with the data, overplotting may be obscuring the true pattern.

Points sitting far from the main cluster may be outliers worth
investigating. It is also important to remember that a relationship
visible in a scatter plot indicates association, not causation. Two
variables can be strongly correlated without one causing the other, and
apparent relationships can sometimes reflect the influence of a third
variable not included in the plot.

\subsection{Numerical Summaries}\label{numerical-summaries}

\subsubsection{Means}\label{means}

\subsubsection{Standard Deviations}\label{standard-deviations}

\subsubsection{Frequency Tables}\label{frequency-tables}




\end{document}
